% !TEX root = task-24.tex
% EGE 2027, task 24. Source IDs and URLs are kept in comments.
% Editable task bank. Append new BankTask blocks at the end of the document.
\documentclass[12pt,a4paper]{article}
\usepackage{physics-textbook}
\usepackage{graphicx,needspace}
\usetikzlibrary{decorations.pathmorphing}
\newif\ifShowSolutions
\ifdefined\TutorEdition
  \ShowSolutionstrue
\else
  \ShowSolutionsfalse % Change to \ShowSolutionstrue to display solutions.
\fi
\newcounter{banktask}
\newsavebox{\BankTaskBox}
\newcommand{\BankTaskContent}[3]{%
  \noindent\textbf{\thebanktask.}\quad #2\par
  \ifShowSolutions
    \Needspace{50mm}% Keep the solution heading with its opening text or diagram.
    #3\par
    \ifdefined\TutorEdition
      \ifcsname TutorSource#1\endcsname
        \nopagebreak\smallskip{\small\csname TutorSource#1\endcsname\par}%
      \fi
    \fi
  \fi
}
\newcommand{\BankTask}[3]{%
  \refstepcounter{banktask}\par\medskip
  \sbox{\BankTaskBox}{\begin{minipage}{\linewidth}%
    \BankTaskContent{#1}{#2}{#3}\end{minipage}}%
  \ifdim\dimexpr\ht\BankTaskBox+\dp\BankTaskBox\relax>\dimexpr\textheight-20mm\relax
    \Needspace{35mm}\BankTaskContent{#1}{#2}{#3}%
  \else
    \Needspace{\dimexpr\ht\BankTaskBox+\dp\BankTaskBox\relax}%
    \noindent\usebox{\BankTaskBox}\par
  \fi
  \medskip
}
\begin{document}


% Source 19871: https://phys-ege.sdamgia.ru/problem?id=19871
\BankTask{19871}{
В комнате при $20$ °С относительная влажность $40$\%. При физической нагрузке через легкие человека приходит $15$ л воздуха за $1$ минуту. Выдыхаемый воздух имеет температуру $34$ °C и влажность $100$\%. Давление насыщенного пара при $20$ °С равно $2{,}34$ кПа, а при $34$ °C составляет $5{,}32$ кПа. Какую массу воды теряет тело человека за $1$ час за счет дыхания? Считать, что выдыхаемый воздух имеет такой же объем, какой проходит через легкие человека. Влажность воздуха в комнате не меняется.
}{
\smallskip\noindent\textbf{Решение}\par
$1$.  Парциальное давление пара при температуре t  =  $20$ °C равно $p=\varphi p_{\,\text{нп}1}.$ Из уравнения Клапейрона  — Менделеева $pV= \frac{m_1}{M} RT_1$ находим массу пара в первом состоянии $m_1= \frac{\varphi p_{\,\text{нп}1}VM}{RT_1} .$ Аналогично для пара во втором состоянии $m_2= \frac{\varphi p_{\,\text{нп}2}VM}{RT_2} ,$ его масса равна $m_2= \frac{p_{\,\text{нп}2}VM}{RT_2} .$
За время $\tau=1\,\text{ч} =60 \,\text{мин}$ человек выдыхает объем водяного пара, равный $V= \frac{\Delta V}{\Delta\tau} \cdot \tau,$ где по условию задачи $\frac{\Delta V}{\Delta\tau} =15 \cdot 10 ^{ - 3 } \,\text{м} ^3 /\,\text{мин}.$ Вычислим объем: V  =  $15$ · $10^{-3}$ · $60$  =  $0{,}9\,\text{м}^{3}$.
$3$.  За $1$ час масса потерянной человеком воды равна:

\[\Delta m=m_2 - m_1= \frac{VM}{R} ( \frac{p_{\,\text{нп}2}}{T_2} - \frac{\varphi p_{\,\text{нп}1}}{T_1} ) = \frac{VM ( p_{\,\text{нп}2}T_1 - \varphi p_{\,\text{нп}1}T_2 )}{RT_1T_2} ,\]

\[\Delta m= \frac{0{,}9 \cdot 18 \cdot 10 ^{ - 3 } ( 5{,}32 \cdot 293 - 2{,}34 \cdot 0{,}4 \cdot 307 ) \cdot 10 ^3}{8{,}31 \cdot 307 \cdot 293} \approx 0{,}028\,\text{кг}=28\,\text{г}.\]

Ответ: $28\,\text{г}$.
}

% Source 7166: https://phys-ege.sdamgia.ru/problem?id=7166
\BankTask{7166}{
\begin{center}\begin{tikzpicture}[task graph,>=stealth]\draw[->](0,0)--(5,0)node[below]{\(T\)};\draw[->](0,0)--(0,3.7)node[left]{\(p\)};\node[below left]at(0,0){\(0\)};\fill(1,1.15)circle(1.25pt)node[left]{\(1\)};\fill(1.85,2.12)circle(1.25pt)node[above]{\(2\)};\fill(3.6,2.12)circle(1.25pt)node[above]{\(3\)};\draw(1,1.15)--(1.85,2.12);\draw[->](1,1.15)--(1.527,1.751);\draw(1.85,2.12)--(3.6,2.12);\draw[->](1.85,2.12)--(2.935,2.120);\draw[thin](0,0)--(1,1.15);\end{tikzpicture}\end{center}
Один моль одноатомного идеального газа совершает процесс $1$–$2$–$3$, график которого показан на рисунке в координатах p–T. Известно, что давление газа p в процессе $1$–$2$ увеличилось в $2$ раза. Какое количество теплоты было сообщено газу в процессе $1$–$2$–$3$, если его температура T в состоянии $1$ равна $300$ К, а в состоянии $3$ равна $900$ К?
}{
\smallskip\noindent\textbf{Решение}\par
Для определения количества теплоты $Q_{123}$ необходимо сложить количества теплоты, сообщенные газу на участках $1$–$2$ и $2$–$3$: $Q_{123}$  =  $Q_{12}$ + $Q_{23}$.
Исходя из приведенного графика, можно сделать вывод, что процесс $1$–$2$ является изохорным. Для него, как следует из уравнения Клапейрона  — Менделеева, $\frac{p_1}{T_1} = \frac{p_2}{T_2} ,$ откуда $\frac{p_2}{p_1} = \frac{T_2}{T_1} =2.$ Следовательно:

\[T_2=T_1 \frac{p_2}{p_1} =2T_1=300 \cdot 2=600\,\text{К}.\]

Работа газа в процессе $1$–$2$ равна нулю, и для него первый закон термодинамики с учетом выражения для внутренней энергии одноатомного идеального газа принимает вид:

\[Q_{12}=\Delta U_{12}= \frac{3}{2} \nu R ( T_2 - T_1 ) = \frac{3}{2} \nu RT_1.\]

Процесс $2$–$3$ является изобарным с давлением p  =  $p_{2}$  =  const, для него первый закон термодинамики принимает вид: $Q=A_{23} + \Delta U_{23},$ где $\Delta U_{23}= \frac{3}{2} \nu R ( T_3 - T_2 )$  — изменение внутренней энергии газа, $A_{23}=p_2 ( V_3 - V_2 )$  — совершенная газом работа. Из уравнения Клапейрона  — Менделеева $pV=\nu RT$ следует, что:

\[Q_{23}= \frac{3}{2} \nu R ( T_3 - T_2 ) + \nu R ( T_3 - T_2 ) = \frac{5}{2} \nu R ( T_3 - T_2 ) .\]

Таким образом:

\[Q_{23}= \frac{5}{2} \nu R ( T_3 - 2T_1 ) .\]

В результате:

\[Q_{123}= \frac{3}{2} \nu RT_1 + \frac{5}{2} \nu R ( T_3 - 2T_1 ) = \frac{3}{2} \cdot 1 \cdot 8{,}31 \cdot 300 + \frac{5}{2} \cdot 1 \cdot 8{,}31 \cdot ( 900 - 600 ) \approx10\,\text{кДж}.\]

Ответ: $10\,\text{кДж}$.
}

% Source 50819: https://phys-ege.sdamgia.ru/problem?id=50819
\BankTask{50819}{
\begin{center}\begin{tikzpicture}[task graph,>=stealth]\draw[->](0,0)--(5,0)node[below]{\(V\)};\draw[->](0,0)--(0,3.7)node[left]{\(p\)};\node[below left]at(0,0){\(0\)};\fill(0.7,2.85)circle(1.25pt)node[above left]{\(1\)};\fill(2.75,2.85)circle(1.25pt)node[above right]{\(2\)};\fill(2.75,0.8)circle(1.25pt)node[right]{\(3\)};\draw(0.7,2.85)--(2.75,2.85);\draw[->](0.7,2.85)--(1.971,2.850);\draw(2.75,2.85)--(2.75,0.8);\draw[->](2.75,2.85)--(2.750,1.579);\draw(2.75,.8)..controls(1.7,.83)and(.77,1.4)..(.7,2.85);\draw[->](1.53,1.23)--(1.3,1.45);\end{tikzpicture}\end{center}
В качестве рабочего тела в тепловой машине используется идеальный одноатомный газ, который совершает циклический процесс, состоящий из изобарного нагревания ($1$ → $2$), изохорного охлаждения ($2$ → $3$) и адиабатного сжатия ($3$ → $1$). Отношение работы $A_{12}$, совершённой газом в изобарном процессе, к работе $A'_{31},$ совершённой над газом при адиабатном сжатии равно $2$. Определите КПД этой тепловой машины.
}{
\smallskip\noindent\textbf{Решение}\par
Проанализируем процессы. В изобарном процессе $1$–$2$ объем увеличивается. По закону Гей  — Люссака $\frac{V}{T} = const$ при увеличении объема увеличивается абсолютная температура газа. 
Работа газа в этом процессе равна $A_{12}=\nu R ( T_2 - T_1 )$ положительна, изменение внутренней энергии одноатомного газа $\Delta U_{12}= \frac{3}{2} \nu R ( T_2 - T_1 )$ положительно. Тогда по первому закону термодинамики $Q_{12}=A_{12} + \Delta U_{12}$ положительно и равно $Q_{12}= \frac{5}{2} \nu R ( T_2 - T_1 ) = \frac{5}{2} A_{12} .$ Значит, газ в этом процессе получает теплоту.
В изохорном процессе $2$–$3$ газ работу не совершает. Согласно закону Шарля $\frac{p}{T} = const$ при уменьшении давления уменьшается абсолютная температура газа. Следовательно, внутренняя энергия газа уменьшается, а по первому закону термодинамики $Q_{23}=\Delta U_{23}$ количество теплоты отрицательно, газ теплоту отдает.
В адиабатном процессе $3$–$1$ газ не получает и не отдает теплоту $Q_{31}=0,$ газ сжимается и поэтому совершает отрицательную работу, модуль которой по условию $|A_{31}|= \frac{A_{12}}{2} .$
Таким образом, работа газа за цикл равна $A=A_{12} + A_{31}= \frac{A_{12}}{2} .$ 
Количество теплоты, полученное газом за цикл $Q=Q_{12}= \frac{5}{2} A_{12} .$
Находим КПД цикла:

\[\eta= \frac{A}{Q} = \frac{2A_{12}}{2 \cdot 5A_{12}} =0{,}2=20\%.\]

Ответ: $20$\%.
}

% Source 38662: https://phys-ege.sdamgia.ru/problem?id=38662
\BankTask{38662}{
Вертикальный цилиндрический сосуд с влажным воздухом влажностью $\varphi$ закрыт невесомым поршнем площадью S, способным скользить без трения. В начальном состоянии поршень находится в равновесии. На поршень медленно насыпают песок. Определите минимальную массу песка m, при которой на стенках сосуда появится роса. Температура в сосуде поддерживается постоянной. Внешнее атмосферное давление равняется $p_{0}$.
}{
\smallskip\noindent\textbf{Решение}\par
В первом состоянии поршень невесомый, поэтому давление влажного воздуха равно атмосферному давлению $p_1=p_0.$ Из формулы влажности $\varphi = \frac{p_{\,\text{п}1}}{p_{\,\text{нп}}} ,$ где $p_{\,\text{п}1}$  — парциальное давление водяного пара в первом состоянии, $p_{\,\text{нп}}$  — давление насыщенного пара при данной температуре.
Во втором состоянии выпала роса, что означает  — пар стал насыщенным. Следовательно,  в сосуде давление водяных паров $p_{\,\text{п}2}$ стало равным $p_{\,\text{нп}}$. 
Поршень в равновесии, следовательно, $p_2 = p_0 + \frac{mg}{S} ,$  где $p_{2}$  — давление пара после насыпания песка.
И для ненасыщенного водяного пара, и для влажного воздуха можно применить закон Бойля-Мариотта при изотермическом процессе:

$\frac{V_2}{V_1} = \frac{p_1}{p_2}$
$\frac{V_2}{V_1} = \frac{p_{\,\text{п}1}}{p_{\,\text{п}2}} ,$

где $V_{1}$ и $V_{2}$ начальный и конечные объемы влажного воздуха в сосуде под поршнем. Из этих соотношений получаем:

\[\frac{p_1}{p_2} = \frac{p_{\,\text{п}1}}{p_{\,\text{п}2}} .\]

 Объединяя эти уравнения, получаем, что масса песка равна

\[m= \frac{p_0S ( 1 - \varphi )}{\varphi g} .\]

Ответ: $m = \frac{p_0 S}{g} \cdot ( \frac{1 - \varphi}{\varphi} ) .$
}

% Source 33046: https://phys-ege.sdamgia.ru/problem?id=33046
\BankTask{33046}{
В вертикальном цилиндре с гладкими стенками под массивным металлическим поршнем находится одноатомный идеальный газ. В начальном состоянии поршень массой M и площадью основания S покоится на высоте h, опираясь на выступы (см. рис. $1$). Давление газа $p_{0}$ равно внешнему атмосферному. Какое количество теплоты Q нужно сообщить газу при медленном его нагревании, чтобы поршень оказался на высоте H (cм. рис. $2$)? Тепловыми потерями пренебречь.

\begin{center}\begin{tikzpicture}[task graph,>=stealth]\draw (0,3.5)--(0,0)--(2.25,0)--(2.25,3.5);\draw(-.1,0)--(2.4,0);\foreach \x in {0,.16,...,2.4}{\draw(\x,0)--++(-.15,-.16);}\fill(0,1.2)--(.23,1.2)--(0,.97)--cycle;\fill(2.25,1.2)--(2.02,1.2)--(2.25,.97)--cycle;\draw[fill=orange!30](0,1.2)rectangle++(2.25,.32);\draw(2.25,1.2)--(3.1,1.2);\draw[<->](2.85,0)--(2.85,1.2)node[midway,right]{\(h\)};\node at(1.05,2.45){\(p_0\)};\node at(1.1,.53){\(p_0\)};\draw(1.8,1.52)--(2.65,2.2)node[right]{\(M,S\)};\draw[->](3.45,1.53)--(4.3,1.53);\end{tikzpicture}\end{center}
Рис. $1$

\begin{center}\begin{tikzpicture}[task graph,>=stealth]\draw (0,3.5)--(0,0)--(2.25,0)--(2.25,3.5);\draw(-.1,0)--(2.4,0);\foreach \x in {0,.16,...,2.4}{\draw(\x,0)--++(-.15,-.16);}\fill(0,1.2)--(.23,1.2)--(0,.97)--cycle;\fill(2.25,1.2)--(2.02,1.2)--(2.25,.97)--cycle;\draw[fill=orange!30](0,2.35)rectangle++(2.25,.32);\draw(2.25,2.35)--(3.1,2.35);\draw[<->](2.85,0)--(2.85,2.35)node[midway,right]{\(H\)};\node at(1.2,3.15){\(p_0\)};\end{tikzpicture}\end{center}
Рис. $2$
}{
\smallskip\noindent\textbf{Решение}\par
Изобразим силы, действующие на поршень в процессе движения.
\begin{center}\begin{tikzpicture}[task graph,>=stealth]\draw[fill=orange!30](-.78,-.25)rectangle(.78,.25);\fill (0,0)circle(1.3pt);\draw[->](0,0)--(0,1.55)node[right]{\(p_2S\)};\draw[->](0,0)--(0,-.77)node[right]{\(Mg\)};\draw[->](0,0)--(0,-1.4)node[right]{\(p_0S\)};\end{tikzpicture}\end{center}
До момента, пока $p_2S < p_0S + Mg,$ поршень будет покоиться, давление газа увеличиваться. При достижении равновесия $p_2S=p_0S + Mg$ поршень начнет двигаться вверх, так как нагревание медленное, из-за чего процесс можно считать изобарным.
По первому закону термодинамики $Q=A + \Delta U.$
Изменение внутренней энергии равно $\Delta U= \frac{3}{2} \nu R\Delta T.$
Работа газа при изобарном расширении:

\[A=p_2\Delta V= ( p_0 + \frac{Mg}{S} ) S ( H - h ) = ( p_0S + Mg ) ( H - h ) .\]

Пусть начальный объем $V_0=Sh,$ конечный объем $V=SH.$
Запишем уравнение Клапейрона-Менделеева для каждого состояния газа:

$p_0V_0=\nu RT_0;$
$p_2V_=\nu RT_2.$

При объединении уравнений, получаем:

\[( p_0S + Mg ) H - p_0Sh=\nu R\Delta T.\]

Тогда изменение внутренней энергии газа равно:

\[\Delta U= \frac{3}{2} ( ( p_0S + Mg ) H - p_0Sh ) .\]

Тогда полученное газом количество теплоты равно:

$Q= \frac{3}{2} ( ( p_0S + Mg ) H - p_0Sh ) + ( p_0S + Mg ) ( H - h ) ;$
$Q= 2{,}5p_0S ( H - h ) + Mg ( 2{,}5H - h ) .$

Ответ: $Q= 2{,}5p_0S ( H - h ) + Mg ( 2{,}5H - h ) .$
}

% Source 32562: https://phys-ege.sdamgia.ru/problem?id=32562
\BankTask{32562}{
В закрытом сосуде объемом V  =  $30$ л находится влажный воздух массой m  =  $45\,\text{г}$ при температуре t  =  $80$ °C и давлении p  =  $160$ кПа. Определите относительную влажность воздуха, находящегося в сосуде. Давление $p_{\,\text{н}}$ насыщенных паров воды при температуре $80$ °С равно $47$ кПа.
}{
\smallskip\noindent\textbf{Решение}\par
Относительная влажность воздуха определяется по формуле $\varphi = \frac{p_{\,\text{п}}}{p_{\,\text{н}}} \cdot 100\%.$
По закону Дальтона давление в сосуде складывается из парциального давления водяного пара $p_{\,\text{п}}$ и давления сухого воздуха $p_{\,\text{возд}}$, то есть $p=p_{\,\text{п}} + p_{\,\text{возд}}.$
Из уравнения Клапейрона  — Менделеева найдем давление сухого воздуха:

$p_{\,\text{возд}}V= \frac{m_{\,\text{возд}}RT}{M_{\,\text{возд}}} ,$ следовательно,  $p_{\,\text{возд}}= \frac{m_{\,\text{возд}}RT}{VM_{\,\text{воз}}} .$

Аналогично находим парциальное давление водяного пара:

$p_{\,\text{п}}V= \frac{m_{\,\text{п}}RT}{M_{\,\text{вод}}} ,$ следовательно, $p_{\,\text{п}}= \frac{m_{\,\text{п}}RT}{VM_{\,\text{вод}}} .$

Учтем, что масса влажного воздуха равна $m=m_{\,\text{п}} + m_{\,\text{возд}}.$ Объединяя уравнения, выразим парциальное давление водяного пара:

\[p_{\,\text{п}}= \frac{pVM_{\,\text{возд}} - mRT}{V ( M_{\,\text{возд}} - M_{\,\text{вод}} )} .\]

Тогда относительная влажность воздуха равна:

\[\varphi = \frac{pVM_{\,\text{возд}} - mRT}{p_{\,\text{н}}V ( M_{\,\text{возд}} - M_{\,\text{вод}} )} \cdot 100\%,\]

 то есть:

\[\varphi = \frac{160 \cdot 10 ^3 \cdot 30 \cdot 10 ^{ - 3 } \cdot 29 \cdot 10 ^{ - 3 } - 45 \cdot 10 ^{ - 3 } \cdot 8{,}31 \cdot 353}{47 \cdot 10 ^3 \cdot 30 \cdot 10 ^{ - 3 } \cdot ( 29 \cdot 10 ^{ - 3 } - 18 \cdot 10 ^{ - 3 } )} \cdot 100\%\approx 46\%.\]

Ответ: $46$\%.
}

% Source 25384: https://phys-ege.sdamgia.ru/problem?id=25384
\BankTask{25384}{
\begin{center}\begin{tikzpicture}[task graph,>=stealth]\draw[->](0,0)--(5,0)node[below]{\(V\)};\draw[->](0,0)--(0,3.7)node[left]{\(p\)};\node[below left]at(0,0){\(0\)};\fill(1,0.85)circle(1.25pt)node[below left]{\(1\)};\fill(1,3.05)circle(1.25pt)node[above left]{\(2\)};\fill(3.85,0.85)circle(1.25pt)node[below right]{\(3\)};\draw(1,0.85)--(1,3.05);\draw[->](1,0.85)--(1.000,2.214);\draw(3.85,0.85)--(1,0.85);\draw[->](3.85,0.85)--(2.083,0.850);\draw(1,3.05)..controls(1.45,1.95)and(2.55,1.58)..(3.85,.85);\draw[->](1.9,1.91)--(2.2,1.72);\end{tikzpicture}\end{center}
$1$ моль одноатомного идеального газа совершает цикл $1$–$2$–$3$–$1$, состоящий из изохоры ($1$–$2$), адиабаты ($2$–$3$) и изобары ($3$–$1$). Абсолютные температуры газа $T_{1}$  =  $400$ K, $T_{2}$  =  $600$ K, $T_{3}$  =  $510$ K. Определите КПД цикла.
}{
\smallskip\noindent\textbf{Решение}\par
КПД циклического процесса равно $\eta = \frac{A}{Q_{\,\text{пол}}} .$ Найдем работу за цикл $A=A_1{,}2 + A_2{,}3 + A_3{,}1.$ Для изохорного процесса $A_1{,}2=0.$ При адиабатном процессе из первого закона термодинамики $Q_2{,}3=A_2{,}3 + \Delta U_2{,}3$ следует, что $A_2{,}3= - \Delta U_2{,}3,$ так как $Q_2{,}3=0,$ а изменение внутренней энергии $\Delta U_2{,}3= \frac{3}{2} \nu R ( T_3 - T_2 ) .$ Для изобарного процесса $A_3{,}1=p\Delta V.$ Учитывая, что из уравнения Клапейрона  — Менделеева $pV=\nu RT,$ то $A_3{,}1=\nu R ( T_1 - T_3 ) .$ Таким образом, работа газа за цикл:

\[A=\nu R ( \frac{3}{2} ( T_2 - T_3 ) + ( T_1 - T_3 ) ) .\]

Из графика следует, что при изохорном процессе газ получал теплоту, при адиабатном процессе количество теплоты равно $0$, при изобарном процессе газ теплоту отдавал. Значит, необходимо найти количество теплоты при изохорном процессе. Из первого закона термодинамики $Q_1{,}2=A_1{,}2 + \Delta U_1{,}2,$ но при этом $A_1{,}2=0,$ а $\Delta U_1{,}2= \frac{3}{2} \nu R ( T_2 - T_1 ) ,$ то от нагревателя получена теплота $Q_{\,\text{пол}}= \frac{3}{2} \nu R ( T_2 - T_1 ) .$
Тогда КПД цикла:

\[\eta = \frac{\frac{3}{2} ( T_2 - T_3 ) + ( T_1 - T_3 )}{\frac{3}{2} ( T_2 - T_1 )} = \frac{1{,}5 \cdot ( 600 - 510 ) + ( 400 - 510 )}{1{,}5 \cdot ( 600 - 400 )} \approx 0{,}08=8\% .\]

Ответ: $8$\%.
}

% Source 43027: https://phys-ege.sdamgia.ru/problem?id=43027
\BankTask{43027}{
Общая масса газонепроницаемой оболочки воздушного шара и его корзины равна $400\,\text{кг}$. Шар заполнен гелием массой $100\,\text{кг}$. Он может поднять в воздух груз массой $225\,\text{кг}$. Какую минимальную массу гелия нужно добавить в оболочку шара, чтобы шар поднял еще одного пассажира массой $50\,\text{кг}$? Считать, что оболочка шара не оказывает сопротивления изменению его объема, воздушных течений в вертикальном направлении нет. Температура и давление гелия внутри шара и воздуха снаружи шара соответственно одинаковы.
}{
\smallskip\noindent\textbf{Решение}\par
Обозначим массу оболочки, гелия и груза $M=m_o + m_{\,\text{г}} + m.$ На шар действуют сила тяжести и выталкивающая сила со стороны воздуха. При поднятии шара подъемная сила равна 

\[F_{\,\text{под}}=F_A - Mg,\]

где выталкивающая сила $F_A=\rho_{\,\text{воз}}gV,$ V - объем шара. Плотность воздуха найдем из уравнения Менделеева-Клапейрона

$p= \frac{\rho_{\,\text{воз}}RT}{M_{\,\text{воз}}} ,$ откуда $\rho_{\,\text{воз}}= \frac{pM_{\,\text{воз}}}{RT} ,$

где p - давление наружного воздуха, T − температура воздуха снаружи шара.
Объем шара найдем из уравнения Менделеева-Клапейрона для гелия:

$pV= \frac{mRT}{M_{\,\text{г}}} ,$ откуда $V= \frac{mRT}{pM_{\,\text{г}}} ,$

где p - давление гелия, равное внешнему давлению, T - температура гелия, равная наружной температуре по условию.
Таким образом, подъемная сила шара равна:

\[F_{\,\text{под}}=g ( \frac{M_{\,\text{воз}}m}{M_{\,\text{г}}} - m_o - m - m_{\,\text{г}} ) =10 \cdot ( \frac{29 \cdot 10 ^{ - 3 } \cdot 100}{4 \cdot 10 ^{ - 3 }} - 400 - 100 - 225 ) =0.\]

Это означает, что масса поднимаемого груза является максимальной.
При добавлении пассажира и минимальной массы гелия в шар подъемная сила не изменится и будет выполняться аналогичное первому случаю соотношение:

\[m_o + m + m_{\,\text{г}} + m_{\,\text{п}}= \frac{M_{\,\text{воз}} ( m + \Delta m )}{M_{\,\text{г}}} .\]

При объединении соотношений для двух случаев получаем:

\[m_{\,\text{п}} + \Delta m= \frac{M_{\,\text{воз}}\Delta m}{M_{\,\text{г}}} ,\]

откуда находим минимальную массу гелия, который нужно добавить в шар:

\[\Delta m= \frac{m_{\,\text{п}}M_{\,\text{г}}}{M_{\,\text{воз}} - M_{\,\text{г}}} = \frac{50 \cdot 4 \cdot 10 ^{ - 3 }}{( 29 - 4 ) \cdot 10 ^{ - 3 }} =8\,\text{кг}.\]

Ответ: $8\,\text{кг}$.
}

% Source 24380: https://phys-ege.sdamgia.ru/problem?id=24380
\BankTask{24380}{
\begin{center}\begin{tikzpicture}[task graph,>=stealth]\draw[->](0,0)--(5,0)node[below]{\(T\)};\draw[->](0,0)--(0,3.7)node[left]{\(p\)};\node[below left]at(0,0){\(0\)};\fill(1,1.4)circle(1.25pt)node[above left]{\(1\)};\fill(2,2.8)circle(1.25pt)node[above left]{\(2\)};\fill(4,2.8)circle(1.25pt)node[above right]{\(3\)};\fill(2,1.4)circle(1.25pt)node[right]{\(4\)};\draw(1,1.4)--(2,2.8);\draw[->](1,1.4)--(1.620,2.268);\draw(2,2.8)--(4,2.8);\draw[->](2,2.8)--(3.240,2.800);\draw(4,2.8)--(2,1.4);\draw[->](4,2.8)--(2.760,1.932);\draw(2,1.4)--(1,1.4);\draw[->](2,1.4)--(1.380,1.400);\draw[thin](0,0)--(1,1.4)(0,0)--(2,1.4)(1,0)--(1,1.4)(2,0)--(2,2.8)(4,0)--(4,2.8);\node[below]at(1,0){\(T_1\)};\node[below]at(2,0){\(T_2=T_4=2T_1\)};\node[below]at(4,0){\(T_3\)};\end{tikzpicture}\end{center}
В тепловом двигателе $1$ моль одноатомного разреженного газа совершает цикл $1$−$2$−$3$−$4$−$1$, показанный на графике в координатах p–T, где p  — давление газа, Т  — абсолютная температура. Температуры в точках $2$ и $4$ равны и превышают температуру в точке $1$ в $2$ раза. Определите КПД цикла.
}{
\smallskip\noindent\textbf{Решение}\par
$1$.  Рассмотрим процессы, происходящие на каждом участке: 
На участке $1$−$2$ происходит изохорное нагревание ($V_{1}$  =  $V_{2}$), при котором по графику температура увеличивается в $2$ раза. Значит, исходя из закона $\frac{p_2}{p_1} = \frac{T_2}{T_1} ,$ давление увеличивается в $2$ раза. 
На участке $2$−$3$ происходит изобарное нагревание, при котором по закону $\frac{V_3}{V_2} = \frac{T_3}{T_2}$ или $\frac{V_3}{V_1} = \frac{T_3}{2T_1} .$
На участке $3$−$4$ происходит изохорное охлаждение, при котором по закону $\frac{p_4}{p_3} = \frac{T_4}{T_3}$ или $\frac{p_4}{2p_1} = \frac{2T_1}{T_3} .$
На участке $4$−$1$ происходит изобарное охлаждение. Тогда в предыдущем уравнении $\frac{p_1}{2p_1} = \frac{2T_1}{T_3} ,$ откуда $T_{3}$  =  $4$T$_{1}$.
$2$.  Найдем работу газа за цикл. На участках $1$−$2$ и $3$−$4$ объемы не меняются, поэтому работа не совершается. Следовательно, работы за цикл равна:

A  =  $A_{2-3}$ + $A_{4-1}$,
$A_2 - 3= \nu R ( T_3 - T_2 ) ;$ $R ( 4T_1 - 2T_1 ) =2\nu RT_1$
$A_4 - 1=\nu R ( T_1 - T_4 ) = - \nu RT_1$
$A=\nu RT_1.$

$3$.  Газ получает теплоту на участках $1$–$2$ и $2$–$3$. Найдем количество теплоты на каждом участке, применив первый закон термодинамики:

$Q_1 - 2=\Delta U_1 - 2= \frac{3}{2} \nu R ( T_2 - T_1 ) =1{,}5\nu RT_1,$
$Q_2 - 3=\Delta U_2 - 3 + A_2 - 3= \frac{3}{2} \nu R ( T_3 - T_2 ) + \nu R ( T_3 - T_2 ) = \frac{5}{2} \nu R ( 4T_1 - 2T_1 ) =5\nu RT_1,$
$Q_{\,\text{пол}}=Q_1 - 2 + Q_2 - 3=6{,}5\nu RT_1.$

$4$.  Найдем КПД циклического процесса:

\[\eta= \frac{A}{Q_{\,\text{пол}}} = \frac{\nu RT_1}{6{,}5\nu RT_1} = \frac{1}{6{,}5} \approx15\%.\]

Ответ: $15$\%.
}

% Source 50827: https://phys-ege.sdamgia.ru/problem?id=50827
\BankTask{50827}{
В калориметре находится $1\,\text{кг}$ льда при температуре −$5$ °C. Какую массу воды, имеющей температуру $20$ °C, нужно добавить в калориметр, чтобы температура его содержимого после установления теплового равновесия оказалась −$2$ °C? Теплообменом с окружающей средой и теплоёмкостью калориметра пренебречь.
}{
\smallskip\noindent\textbf{Решение}\par
В результате теплообмена происходят процессы нагревания льда от $t_{1}$ до установившейся температуры t, охлаждение воды от $t_{2}$ до $t_{\,\text{пл}}=0 ^\circ \,\text{С},$ кристаллизация воды и охлаждение образовавшегося льда от $t_{\,\text{пл}}$ до установившейся температуры t. Для каждого из этих процессов запишем количество теплоты: 

$Q_1=m_1\,\text{с}_{\,\text{л}} ( t - t_1 ) ;$
$Q_2=m_2c_{\,\text{в}} ( t_{\,\text{пл}} - t_2 ) = - m_2c_{\,\text{в}}t_2;$
$Q_3= - \lambda m_2;$
$Q_4=m_2c_{\,\text{л}} ( t - t_{\,\text{пл}} ) =m_2c_{\,\text{л}}t.$

Так как при теплообмене отсутствует теплообмен с окружающей средой, то уравнение теплового баланса имеет вид:

$Q_1 + Q_2 + Q_3 + Q_4=0,$
$m_1\,\text{с}_{\,\text{л}} ( t - t_1 ) - m_2c_{\,\text{в}}t_2 - \lambda m_2 + m_2c_{\,\text{л}}t=0,$

откуда находим массу воды:

\[m_2= \frac{m_1\,\text{с}_{\,\text{л}} ( t - t_1 )}{c_{\,\text{в}}t_2 + \lambda - c_{\,\text{л}}t} = \frac{1 \cdot 2100 \cdot ( - 2 + 5 )}{4200 \cdot 20 + 330000 + 2100 \cdot 2} \approx 0{,}015\,\text{кг}=15\,\text{г}.\]

Ответ: $15\,\text{г}$.
}

% Source 30001: https://phys-ege.sdamgia.ru/problem?id=30001
\BankTask{30001}{
В запаянной с одного конца трубке находится влажный воздух, отделенный от атмосферы столбиком ртути длиной l  =  $76\,\text{мм}$. Когда трубка лежит горизонтально, относительная влажность воздуха $\varphi_1$ в ней равна $80$\%. Какой станет относительная влажность этого воздуха $\varphi_2,$ если трубку поставить вертикально, открытым концом вниз? Атмосферное давление равно $760\,\text{мм}$ рт. ст. Температуру считать постоянной.
}{
\smallskip\noindent\textbf{Решение}\par
$1$.  Относительная влажность воздуха в трубке в первом случае: $\varphi_1= \frac{p_1}{p_{\,\text{н}}} ,$ где $p_{1}$  — давление водяных паров в трубке, $p_{\,\text{н}}$  — давление насыщенных водяных паров при той же температуре. Для второго случая относительная влажность воздуха в трубке: $\varphi_2= \frac{p_2}{p_{\,\text{н}}} ,$ где $p_{2}$  — парциальное давление водяных паров в трубке.

\begin{center}\begin{tikzpicture}[task graph,>=stealth]\draw(.35,0)arc(-90:-270:.35)--(4.4,.7);\draw(.35,0)--(4.4,0);\draw(4.4,.35)ellipse(.16 and .43);\draw[fill=orange!30](1.62,0)rectangle(3.05,.7);\draw(0,-.1)--(0,-1.05)(1.62,0)--(1.62,-1.05)(3.05,0)--(3.05,-1.05);\draw[<->](0,-.88)--(1.62,-.88)node[midway,above]{\(l_1\)};\draw[<->](1.62,-.88)--(3.05,-.88)node[midway,above]{\(l\)};\node at(3.96,.95){\(p_0\)};\end{tikzpicture}\end{center}

\begin{center}\begin{tikzpicture}[task graph,>=stealth]\draw(0,0)--(0,3.55)arc(180:0:.33)--(.66,0);\draw(.33,0)ellipse(.44 and .15);\draw[fill=orange!30](0,1.26)rectangle(.66,2.48);\draw(.66,1.26)--(1.75,1.26)(.66,2.48)--(1.75,2.48)(.33,3.88)--(1.75,3.88);\draw[<->](1.59,1.26)--(1.59,2.48)node[midway,left]{\(l\)};\draw[<->](1.59,2.48)--(1.59,3.88)node[midway,left]{\(l_2\)};\node at(-.46,.4){\(p_0\)};\end{tikzpicture}\end{center}

$2$.  Давление влажного воздуха в первом случае равно атмосферному давлению: $p_1 \,\text{вл}=p_0,$ а во втором случае  — $p_2\,\text{вл}=p_0 - \rho g l,$ где \ensuremath{\rho} — плотность ртути, а l  — длина столбика ртути. Заметим, что $p_0=\rho g H,$ где H  =  $760\,\text{мм}$. 
$3$.  Объем влажного воздуха при изменении положения трубки изменился так, как изменились длины столбиков воздуха, запертых ртутью. Считаем влажный воздух идеальным газом. Поэтому по закону Бойля  — Мариотта можем записать:

\[\frac{V_2}{V_1} = \frac{l_2}{l_1} = \frac{p_0}{p_0 - \rho g l} .\]

$4$.  Из уравнения Клапейрона  — Менделеева следует, что: 

$p= \frac{\nu R T}{V} ,$ $p_{\,\text{вл}}= \frac{( \nu + \nu_{\,\text{св}} ) R T}{V} ,$

 где p  — парциальное давление водяного пара, $p_{\,\text{вл}}$  — давление влажного воздуха, $\nu$ — количество моль водяного пара в трубке, $\nu_{\,\text{св}}$  — количество моль сухого воздуха в трубке. Отсюда следует, что:

\[\frac{p}{p_{\,\text{вл}}} = \frac{\nu}{\nu + \nu_{\,\text{св}}} = const.\]

Поэтому $\frac{p_1}{p_2} = \frac{p_1 \,\text{вл}}{p_2 \,\text{вл}} ,$ и для изотермического процесса для парциальных давлений водяного пара имеем: 

\[\frac{p_1}{p_2} = \frac{V_2}{V_1} = \frac{p_0}{p_0 - \rho g l} .\]

$5$.  После объединения записанных выше уравнений получим выражение для относительной влажности воздуха во втором случае: 

\[\varphi_2=\varphi_1 \cdot \frac{p_0 - \rho g l}{p_0} =\varphi_1 \cdot \frac{\rho g H - \rho g l}{\rho g H} =\varphi_1 \frac{H - l}{H} =80 \cdot \frac{760 - 76}{760} =72 \%.\]

Ответ: $\varphi_2=72 \%.$
}

% Source 2988: https://phys-ege.sdamgia.ru/problem?id=2988
\BankTask{2988}{
В калориметре находился лед при температуре $t_1= - 5 ^\circ\,\text{С}.$ Какой была масса $m_1$ льда, если после добавления в калориметр $m_2=4\,\text{кг}$ воды, имеющей температуру $t_2=20 ^\circ\,\text{С},$ и установления теплового равновесия температура содержимого калориметра оказалась равной $t=0 ^\circ\,\text{С},$ причем в калориметре была только вода?
}{
\smallskip\noindent\textbf{Решение}\par
Количество теплоты, полученное при нагреве льда, находящегося в калориметре, до температуры $t=0 ^\circ C$: $Q_1=c_1m_1 ( 0 - t_1 )$ ($1$).
Количество теплоты, полученное льдом при его таянии при $0 ^\circ\,\text{С}:$ $Q_2=\lambda m_1$ ($2$).
Количество теплоты, отданное водой при охлаждении ее до $0 ^\circ\,\text{С}:$ $Q=c_2m_2 ( t_2 - 0 )$ ($3$). 
Уравнение теплового баланса: $Q=Q_1 + Q_2$ ($4$). 
Объединяя ($1$)–($4$), получаем: 

\[m_1= \frac{m_2}{c} _2 ( t_2 - 0 ) c_1 ( 0 - t_1 ) + \lambda = \frac{4 \cdot 4200 \cdot 20}{2100 \cdot 5 + 3{,}3 \cdot 10 ^{5}} \approx 1\,\text{кг}.\]

Ответ: $m_1\approx 1\,\text{кг}.$
}

% Source 50825: https://phys-ege.sdamgia.ru/problem?id=50825
\BankTask{50825}{
Сосуд объёмом V  =  $2$ л разделён на две части. В одной части объёмом $V_{1}$ находятся вода и насыщенный водяной пар. Масса воды в этой части в $2$ раза меньше массы пара. После удаления перегородки пар заполняет весь сосуд. Температура во время процесса остаётся постоянной. После установления равновесия относительная влажность в сосуде составляет $30$\%. Определите объём $V_{2}$.
}{
\smallskip\noindent\textbf{Решение}\par
Отметим, что в сосуде первоначально находился насыщенный пар, следовательно, относительная влажность воздуха была $100$\%. Так как температура во время процесса не изменялась, то плотность насыщенного пара оставалась неизменной и равной первоначальной.
В сосуде объемом $V_{1}$ находились вода и пар. Следовательно, общая масса содержимого равнялась 

\[m=m_{\,\text{п}} + m_{\,\text{в}}=m_{\,\text{п}} + \frac{m_{\,\text{п}}}{2} = \frac{3}{2} m_{\,\text{п}}.\]

Плотность пара была равна плотности насыщенного пара $\rho_{\,\text{н}}.$ Учитывая, что масса пара $m_{\,\text{п}}=\rho_{\,\text{н}}V_1,$ получаем, что масса содержимого в первом сосуде равнялась $m= \frac{3}{2} \rho_{\,\text{н}}V_1.$
После удаления перегородки масса содержимого не изменилась. Так как пар стал ненасыщенным, то вся вода испарилась, в сосуде был пар с плотностью $\rho,$ которая стала равной $\rho= \varphi \rho_{\,\text{н}}.$ При этом сохранившаяся масса пара равна $m=\rho V= \varphi \rho_{\,\text{н}}V.$
Учитывая неизменность массы, получаем равенство $\frac{3}{2} \rho_{\,\text{н}}V_1=\varphi \rho_{\,\text{н}}V,$ откуда находим объем первой части сосуда:

\[V_1= \frac{2 \varphi V}{3} .\]

Следовательно, объем второй части сосуда

\[V_2=V - V_1=V - \frac{2 \varphi V}{3} = \frac{3 - 2 \varphi}{3} V= \frac{3 - 2 \cdot 0{,}3}{3} \cdot 2 =1{,}6\,\text{л}.\]

Ответ: $1{,}6$ л.
}

% Source 33071: https://phys-ege.sdamgia.ru/problem?id=33071
\BankTask{33071}{
\begin{center}\begin{tikzpicture}[task graph,>=stealth]\draw(0,.8)--(0,2.7);\foreach \y in {.85,1.1,...,2.6}{\draw(0,\y)--++(-.22,-.21);}\draw(.21,2.09)circle(.19);\draw(.21,2.09)circle(.07);\draw(.21,2.01)rectangle(5,2.17);\node[above]at(.28,2.24){\(A\)};\node[above]at(1.86,2.2){\(B\)};\draw(1.77,2.01)rectangle(1.93,1.25);\draw(1.45,1.03)rectangle(2.23,1.25);\draw[line width=1.4pt](.53,0)--(.53,1.03)--(3.04,1.03)--(3.04,0);\draw(5,2.01)--(5,1.42);\draw[fill=orange!30](4.76,.86)rectangle(5.24,1.42);\end{tikzpicture}\end{center}
В цилиндр объемом $0{,}5\,\text{м}^{3}$ насосом закачивается воздух со скоростью $0{,}002\,\text{кг}$/c. В верхнем торце цилиндра есть отверстие, закрытое предохранительным клапаном. Клапан удерживается в закрытом состоянии стержнем, который может свободно поворачиваться вокруг оси в точке А (см. рис.). К свободному концу стержня подвешен груз массой $2\,\text{кг}$. Клапан открывается через $580\,\text{с}$ работы насоса, если в начальный момент времени давление воздуха в цилиндре было равно атмосферному. Площадь закрытого клапаном отверстия $5 \cdot 10 ^{ - 4 } \,\text{м} ^2 ,$ расстояние АB равно $0{,}1\,\text{м}$. Температура воздуха в цилиндре и снаружи не меняется и равна $300$ K. Определите длину стержня, если его считать невесомым.
}{
\smallskip\noindent\textbf{Решение}\par
Клапан откроется, когда избыточная сила F давления воздуха на клапан изнутри цилиндра сравняется с силой давления стержня на этот клапан. Если превышение давления воздуха в цилиндре над атмосферным $\Delta p,$ а площадь клапана S, то:

\[F=S \cdot \Delta p.\]

Сила действия стержня на клапан равна $mg \cdot \frac{L}{l} ,$ где m, L и l соответственно масса груза, длина стержня и длина его участка АВ. Итак, должно выполняться условие:

\[S \cdot \Delta p > \,\text{или} \,\text{равно} mg \cdot \frac{L}{l} .\]

Дополнительное давление воздуха определяется увеличением массы $\Delta m_{\,\text{в}}$ воздуха в цилиндре. Согласно уравнению Клапейрона  — Менделеева:

\[\Delta p= \frac{\Delta m_{\,\text{в}}}{MV} RT,\]

где M  — молярная масса воздуха. Поэтому условие открытия клапана имеет вид:

\[\frac{S\Delta m_{\,\text{в}}}{MV} RT > \,\text{или} \,\text{равно} mg \cdot \frac{L}{l} ,\]

 или 

\[l > \,\text{или} \,\text{равно} \frac{mgMVL}{SRT\Delta m_{\,\text{в}}} .\]

Если насос закачивает каждую секунду $\omega$ кг воздуха, то за время t он закачает в цилиндр массу $\Delta m_{\,\text{в}}=\omega t.$ Следовательно, клапан откроется в момент, когда выполнится равенство, из которого найдем длину стержня:

\[L= \frac{SRT\omega tl}{mgMV} = \frac{5 \cdot 10 ^{ - 4 } \cdot 8{,}31 \cdot 300 \cdot 0{,}002 \cdot 580 \cdot 0{,}1}{2 \cdot 10 \cdot 29 \cdot 10 ^{ - 3 } \cdot 0{,}5} \approx 0{,}5\,\text{м} .\]

Ответ: $0{,}5\,\text{м}$.
}

% Source 43206: https://phys-ege.sdamgia.ru/problem?id=43206
\BankTask{43206}{
\begin{center}\begin{tikzpicture}[task graph,>=stealth]\draw[->](0,0)--(5,0)node[below]{\(V\)};\draw[->](0,0)--(0,3.7)node[left]{\(p\)};\node[below left]at(0,0){\(0\)};\fill(0.7,3)circle(1.25pt)node[above]{\(1\)};\fill(3.15,3)circle(1.25pt)node[above right]{\(2\)};\fill(3.15,0.47)circle(1.25pt)node[right]{\(3\)};\draw(.7,3)--(3.15,3)--(3.15,.47)..controls(1.38,.55)and(.82,1.5)..(.7,3);\end{tikzpicture}\end{center}
Один моль одноатомного идеального газа совершает цикл $1$–$2$–$3$–$1$, состоящий из изобары ($1$–$2$), изохоры ($2$–$3$) и адиабаты ($3$–$1$). Абсолютные температуры газа в состояниях $1$, $2$ и $3$ равны $400$ К, $800$ К и $252$ К соответственно. Определите КПД теплового двигателя, работающего по этом циклу.
}{
\smallskip\noindent\textbf{Решение}\par
На участке $1$−$2$ при изобарном расширении температура газа также увеличивается согласно закону $\frac{V}{T} =const.$ При этом газ совершает положительную работу $A_{12}=\nu R ( T_2 - T_1 ) ,$ его внутренняя энергия увеличивается $\Delta U_{12}= \frac{3}{2} \nu R ( T_2 - T_1 ) .$ Согласно первому закону термодинамики газ получает от нагревателя количество теплоты

\[Q_{12}=A_{12} + \Delta U_{12}= \frac{5}{2} \nu R ( T_2 - T_1 ) .\]

На участке $2$−$3$ давление изохорно уменьшается. Согласно закону $\frac{p}{T} =const$ температура газа уменьшается. Тогда газ работу не совершает, а его внутренняя энергия уменьшается $\Delta U_{23}= \frac{3}{2} \nu R ( T_3 - T_2 ) .$ По первому закону термодинамики газ отдает холодильнику количество теплоты

\[|Q_{23}|=A_{23} + \Delta U_{23}= \frac{3}{2} \nu R|T_3 - T_2|= \frac{3}{2} \nu R ( T_2 - T_3 ) .\]

Процесс $3$−$1$ является адиабатным, поэтому $Q_{31}=0.$
Тогда КПД теплового двигателя:

\[\eta=1 - \frac{Q_{\,\text{х}}}{Q_{\,\text{н}}} =1 - \frac{|Q_{23}|}{Q_{12}} =1 - \frac{3 ( T_2 - T_3 )}{5 ( T_2 - T_1 )} =1 - \frac{3 \cdot ( 800 - 252 )}{5 \cdot ( 800 - 400 )} =0{,}178=17{,}8\%.\]

Ответ: $17{,}8$\%.
}

% Source 9042: https://phys-ege.sdamgia.ru/problem?id=9042
\BankTask{9042}{
В вакууме закреплен горизонтальный цилиндр (см. рис.). В цилиндре находится гелий, запертый поршнем. 
\begin{center}\begin{tikzpicture}[task graph,>=stealth]\draw[fill=gray!38](0,0)rectangle(5.7,1.85);\draw[fill=white](.8,.47)rectangle(5.7,1.38);\draw[fill=orange!35](.8,.47)rectangle(3.33,1.38);\fill(3.33,.47)rectangle(3.56,1.38);\draw[fill=gray!38](3.56,.47)--(4.13,.47)--(3.56,.72)--cycle;\draw[fill=gray!38](3.56,1.38)--(4.13,1.38)--(3.56,1.12)--cycle;\draw[->](6.6,.92)--(5.72,.92);\end{tikzpicture}\end{center}
Поршень массой $90\,\text{г}$ удерживается упорами и может скользить влево вдоль стенок цилиндра без трения. В поршень попадает пуля массой $10\,\text{г}$, летящая горизонтально со скоростью $400\,\text{м/с}$, и застревает в нем. Температура гелия в момент остановки поршня в крайнем левом положении возрастает на $64$ К. Чему равно количество вещества гелия в цилиндре? Считать, что за время движения поршня газ не успевает обменяться теплом с цилиндром и поршнем.
}{
\smallskip\noindent\textbf{Решение}\par
$1$.  Запишем закон сохранения импульса:

\[m v _1= ( m + M ) v _2,\]

где m и M  — масса пули и поршня соответственно, $v _1$  — скорость пули, $v _2$  — скорость поршня с застрявшей пулей.
Поршень будет двигаться со скоростью:

\[v _2= \frac{m v _1}{m + M} .\]

$2$.  Поршень с пулей будет обладать кинетической энергией:

\[E_{\,\text{кин}}= \frac{( m + M ) v _2 ^2}{2} = \frac{m ^2 v _1 ^2}{2 ( m + M )} .\]

Изменение кинетической энергии поршня будет равно работе газа $A=\Delta E_{\,\text{к}}=E_{\,\text{к}2} - E_{\,\text{к}1}= - \frac{m ^2 v _1 ^2}{2 ( m + M )} .$ Вследствие остановки поршня $E_{\,\text{к}2}=0.$
$3$.  Поскольку газ не успевает обменяться теплом с цилиндром и поршнем, то сжатие газа будет являться адиабатичным процессом. По первому началу термодинамики $- A=\Delta U.$ Вследствие этого процесса вся механическая энергия движения поршня с пулей пойдет на нагрев газа, поэтому:

\[\Delta U= \frac{3}{2} \nu R \Delta T= \frac{m ^2 v _1 ^2}{2 ( m + M )} .\]

$4$.  Найдем отсюда количество вещества в цилиндре:

\[\nu= \frac{m ^2 v _1 ^2}{3 ( m + M ) R \Delta T} = \frac{0{,}01 ^2 \cdot 400 ^2}{3 \cdot ( 0{,}01 + 0{,}09 ) \cdot 8{,}31 \cdot 64} \approx 0{,}1\,\text{моль}.\]

Ответ: $\nu\approx 0{,}1\,\text{моль}.$
}

% Source 29749: https://phys-ege.sdamgia.ru/problem?id=29749
\BankTask{29749}{
Закрытый вертикальный цилиндрический сосуд, заполненный идеальным газом, разделён на две части тяжёлым поршнем, способным скользить без трения. В начальном равновесном состоянии в верхней и нижней частях сосуда находилось по $\nu$  =  $1$ моль газа, а отношение объёмов верхней и нижней частей сосуда было равно $2$. После того как из верхней части сосуда полностью откачали газ, через длительный промежуток времени установилось новое состояние равновесия. Найдите отношение объёмов верхней и нижней частей сосуда после откачки газа. Температура газа T в обеих частях сосуда всё время поддерживалась одинаковой и постоянной.
}{
\smallskip\noindent\textbf{Решение}\par
$1$.  Вначале на поршень действовали: сила давления газа в верхней части цилиндра $F_1=p_1 S,$ направленная вниз, сила давления газа в нижней части цилиндра $F_2=p_2 S,$ направленная вверх, сила тяжести поршня mg, направленная вниз, где $p_{1}$  и $p_{2}$  — давления газа в верхней и нижней частях цилиндра, S  — площадь поршня, m  — масса поршня. Поскольку поршень находится в равновесии, то $p_2 S=p_1 S + m g.$ Учитывая, что давление в верхней части меньше, чем в нижней, делаем вывод, что объем газа в верхней части больше, чем объем газа в нижней части цилиндра. 
Обозначим объем нижней части цилиндра V, объем верхней части цилиндра 2V.
Используя уравнение Клапейрона  — Менделеева $p V= \nu R T,$ выразим давления каждой части газа: $p_1= \frac{\nu R T}{2 V}$ и $p_2= \frac{\nu RT}{V} .$ Тогда сила тяжести, действующая на поршень $m g= \frac{\nu RTS}{2 V} .$ 
$2$.  Когда из верхней части цилиндра удалили весь газ, поршень поднимался до тех пор, пока сила тяжести не уравновесилась силой давления газа, находящегося в нижней части цилиндра, то есть $m g=p S.$ С учетом выраженного значения силы тяжести и применяя уравнение Клапейрона  — Менделеева для этого случая, получаем объем газа в нижней части цилиндра:

\[V_2= \frac{\nu R T}{p} = \frac{\nu R T \cdot 2 V}{\nu R T} =2 V.\]

Таким образом, в нижней части объем газа будет в $2$ раза больше, чем в верхней. Значит, отношение объема верхней части цилиндра к объему нижней части станет равным $0{,}5$. 

Ответ: $0{,}5$.
}

% Source 9248: https://phys-ege.sdamgia.ru/problem?id=9248
\BankTask{9248}{
Мальчик решил подняться в воздух на воздушных шарах с гелием. Известно, что мальчик весит $40\,\text{кг}$, а на улице нормальное атмосферное давление и температура $27$ °C. Учитывая, что объем одного шара составляет $10$ литров, найдите, сколько потребуется шаров для такого путешествия. Массой оболочки шаров и объемом мальчика пренебречь.
}{
\smallskip\noindent\textbf{Решение}\par
На мальчика действуют силы: сила Архимеда $F_A=\rho_{\mathrm{He}} ( nV ) g,$ сила тяжести, действующая на человека, mg и сила тяжести, действующая на шарики, $m_{\mathrm{He}}g=\rho_{\mathrm{He}}g ( nV ) .$
Мальчик будет подниматься за счет силы Архимеда, действующей на шары:

$mg + m_{\mathrm{He}}g=F_A.$
$mg + \rho_{\mathrm{He}} ( nV ) g=g\rho_{\,\text{в}} ( nV ) ,$

откуда количество шариков:

\[n= \frac{m}{V ( \rho_{\,\text{в}} - \rho_{\mathrm{He}} )} ,\]

где $\rho_{\,\text{в}}$  — плотность воздуха, $\rho_{\mathrm{He}}$  — плотность гелия.
Из уравнения Менделеева  — Клапейрона можем найти плотность газа при данной температуре и давлении:

$p= \frac{\rho}{M} RT ,$
$\rho= \frac{pM}{RT} .$

Найдем необходимое количество шаров, учитывая, что гелий в шариках находится при той же температуре и давлении, что и воздух вокруг:

\[n= \frac{m}{V ( \rho_{\,\text{в}} - \rho_{\mathrm{He}} )} = \frac{mRT}{Vp_0 ( M_{\,\text{в}} - M_{\mathrm{He}} )} = \frac{40 \cdot 8{,}31 \cdot 300}{10 \cdot 10 ^{ - 3 } \cdot 10 ^{ 5 } \cdot ( 0{,}029 - 0{,}004 )} \approx 4000.\]

Ответ: $4000$.
}

% Source 29750: https://phys-ege.sdamgia.ru/problem?id=29750
\BankTask{29750}{
\begin{center}\begin{tikzpicture}[task graph,>=stealth]\draw(0,0)--(0,1.75);\draw(.83,.45)--(5.35,.45)--(5.35,1.24)--(.83,1.24);\draw[fill=orange!30](2.45,.45)rectangle(3.13,1.24);\end{tikzpicture}\end{center} Горизонтально закрепленная пробирка со столбиком ртути, длиной $1\,\text{см}$, вращается с угловой скоростью $10$ c$^{-1}$. Во сколько раз нужно увеличить температуру внутри пробирки, чтобы столбик ртути не сдвинулся при увеличении угловой скорости в $4$ раза. Начальная температура ртути $0$ °C. Расстояние от оси вращения до центра массы ртути $20\,\text{см}$. Давление снаружи пробирки считать атмосферным.
}{
\smallskip\noindent\textbf{Решение}\par
\begin{center}\begin{tikzpicture}[task graph,>=stealth]\draw(0,0)--(0,1.75);\draw(.83,.45)--(5.35,.45)--(5.35,1.24)--(.83,1.24);\draw[fill=orange!30](2.45,.45)rectangle(3.13,1.24);\draw[->](.76,.85)--(.11,.85)node[midway,above]{\(\vec a\)};\draw[->](2.45,.85)--(1.4,.85)node[midway,above]{\(\vec F_1\)};\draw[->](3.13,.85)--(4.3,.85)node[midway,above]{\(\vec F_2\)};\end{tikzpicture}\end{center} Горизонтальная пробирка имеет столбик ртути, на который оказывается давление снаружи  — атмосферное, внутри  — давление запертого ртутью воздуха. Тогда на этот столбик ртути действуют силы $F_1=p_1 S$ и $F_2=p_0 S.$ 
$2$.  По второму закону Ньютона $m a=F_1 - F_2.$ При этом центростремительное ускорение $a=\omega ^2 R,$ масса ртути $m=\rho V,$ объем ртути $V=S l.$ 
Тогда, объединив уравнения, найдем давление воздуха, запертого столбиком ртути: 

\[p_1=p_0 + \rho l \omega ^2 R=10 ^{ 5 } + 13{,}6 \cdot 10 ^3 \cdot 0{,}01 \cdot 10 ^2 \cdot 0{,}2=102{,}72 \,\text{кПа}.\]

$3$.  Учитывая, что столбик ртути должен остаться на месте, делаем вывод, что объем воздуха, запертого ртутью, не изменился, следовательно, процесс изохорный, для которого $\frac{p_1}{T_1} = \frac{p_2}{T_2} ,$ где $p_{2}$  — давление столбика воздуха после нагревания. 
$4$.  Объединим формулы, учтем, что угловая скорость увеличилась в $4$ раза, для каждого случая имеем уравнения: 

$p_1 - p_0=\rho l \omega ^2 R;$ $\frac{p_1 T_2}{T_1} - p_0=16 \rho l \omega ^2 R.$

Решая данную систему уравнений, находим отношение конечной и начальной температур воздуха, запертого ртутью в пробирке:

\[\frac{T_2}{T_1} =16 - \frac{15 p_0}{p_1} =16 - \frac{15 \cdot 100}{102{,}72} \approx 1{,}4.\]

Ответ: увеличить в $1{,}4$ раза.
}

% Source 19877: https://phys-ege.sdamgia.ru/problem?id=19877
\BankTask{19877}{
\begin{center}\begin{tikzpicture}[task graph,>=stealth]\foreach \x/\y in {0/1.68,3.4/1.48}{\fill[orange!28](\x,0)rectangle++(2.4,\y);\draw[fill=gray!35](\x,\y)rectangle++(2.4,.25);\draw(\x,2.8)--(\x,0)--++(2.4,0)--++(0,2.8);\draw(\x-.2,0)--++(2.8,0);\foreach \a in {0,.18,...,2.52}{\draw(\x+\a,0)--++(-.16,-.15);}\fill(\x+1.2,\y+.125)circle(1.2pt);\draw(\x+1.2,\y+.125)--++(-.42,.42)node[above left]{\(m\)};}\end{tikzpicture}\end{center} В вертикальном цилиндрическом сосуде с гладкими стенками под поршнем массой m  =  $10\,\text{кг}$ и площадью поперечного сечения S  =  $50\,\text{см}^{2}$ находится разреженный газ (см. рис.). При движении сосуда по вертикали с ускорением, направленным вверх и равным по модулю a  =  $1\,\text{м}$/c$^{2}$, высота столба газа под поршнем постоянна и на $5$ \% меньше, чем в покоящемся сосуде. Считая температуру газа под поршнем неизменной, а наружное давление постоянным, найдите это наружное давление. Масса газа под поршнем постоянна.
}{
\smallskip\noindent\textbf{Решение}\par
Газ разреженный, поэтому его можно считать идеальным. Запишем условия равновесия поршня для каждого случая:
I. $p_1=p_{\,\text{н}} + p_{\,\text{п}1}=p_{\,\text{н}} + \frac{mg}{S} .$
II. $p_2=p_{\,\text{н}} + p_{\,\text{п}2}=p_{\,\text{н}} + \frac{mg + ma}{S} .$
Запишем отношения объемов до и после движения сосудов: $\frac{V_2}{V_1} = \frac{h_2}{h_1} =0{,}95.$ По уравнению Менделеева-Клапейрона: $pV=\nu RT,$ $m_{\,\text{газа}}=const,$ $T=const ,$ следовательно, $pV=const,$ или $p_1V_1=p_2V_2,$ откуда $\frac{p_2}{p_1} = \frac{V_1}{V_2} = \frac{1}{0{,}95} .$ Таким образом, $p_1=0{,}95p_2.$ Откуда получаем:

$p_{\,\text{н}} + \frac{mg}{S} =0{,}95 ( p_{\,\text{н}} + \frac{mg + ma}{S} ) ;$
$p_{\,\text{н}} - 0{,}95p_{\,\text{н}}= \frac{( mg + ma ) \cdot 0{,}95}{S} - \frac{mg}{S} ;$
$0{,}05p_{\,\text{н}}= \frac{m}{S} ( 0{,}95 a - 0{,}05g ) ;$
$p_{\,\text{н}}= \frac{10}{0{,}05 \cdot 5 \cdot 10 ^{ - 3 }} \cdot ( 0{,}95 \cdot 1 - 0{,}05 \cdot 10 ) =1{,}8 \cdot 10 ^{4} \,\text{Па}.$

Ответ: $p_{\,\text{н}}$  =  $1{,}8$ · $10^{4}$ Па.
}

% Source 35344: https://phys-ege.sdamgia.ru/problem?id=35344
\BankTask{35344}{
\begin{center}\begin{tikzpicture}[task graph,>=stealth]\draw[gray!65,step=.8] (0,0) grid (4,3.2);\draw[->](0,0)--(5,0)node[below]{\(V\)};\draw[->](0,0)--(0,3.7)node[left]{\(p\)};\node[below left]at(0,0){\(0\)};\fill(0.8,0.8)circle(1.25pt)node[below left]{\(1\)};\fill(0.8,1.6)circle(1.25pt)node[above left]{\(2\)};\fill(3.2,2.4)circle(1.25pt)node[above right]{\(3\)};\draw(0.8,0.8)--(0.8,1.6);\draw[->](0.8,0.8)--(0.800,1.296);\draw(0.8,1.6)--(3.2,2.4);\draw[->](0.8,1.6)--(2.288,2.096);\draw(3.2,2.4)--(0.8,0.8);\draw[->](3.2,2.4)--(1.712,1.408);\node[left]at(0,.8){\(p_0\)};\node[left]at(0,2.4){\(3p_0\)};\node[below]at(.8,0){\(V_0\)};\node[below]at(3.2,0){\(4V_0\)};\end{tikzpicture}\end{center}
В цикле, показанном на pV-диаграмме, v  =  $4$ моль разреженного гелия получает от нагревателя количество теплоты $Q_{\,\text{нагр}}$  =  $120\,\text{кДж}$. Найдите температуру $T_{2}$ гелия в состоянии $2$.
}{
\smallskip\noindent\textbf{Решение}\par
Согласно графику цикла гелий получает положительное количество теплоты от нагревателя на участках $1$−$2$ и $2$−$3$. При этом процесс $1$–$2$ является изохорным и газ работы не совершает. 
По первому закону термодинамики для процесса $1$–$2$ $Q_1{,}2=\Delta U_1{,}2,$ изменение внутренней энергии в этом процессе:

\[\Delta U_1{,}2=U_2 - U_1= \frac{3}{2} p_0V_0.\]

Для процесса $2$–$3$ по первому закону термодинамики $Q_2{,}3=\Delta U_2{,}3 + A_2{,}3.$ Изменение внутренней энергии для данного процесса:

\[\Delta U_2{,}3=U_3 - U_2= \frac{3}{2} ( 12p_0V_0 - 2p_0V_0 ) =15p_0V_0.\]

Работу газа для процесса $2$–$3$ найдем графически как площадь фигуры под графиком:

\[A_2{,}3= \frac{p_2 + p_3}{2} ( V_3 - V_2 ) = \frac{2p_0 + 3p_0}{2} ( 4V_0 - V_0 ) = \frac{15}{2} p_0V_0.\]

Тогда полученное количество теплоты равно: 

\[Q_{\max p} =Q_{12} + Q_{23} = \frac{3}{2} p_0V_0 + 15p_0V_0 + \frac{15}{2} p_0V_0 =24 p_0 V_0 .\]

Согласно графику цикла $v R T_2=2 p_0 V_0,$ откуда: 

\[T_2= \frac{2 p_0 V_0}{v R} = \frac{Q_{\,\text{нагр}}}{12 v R} = \frac{120 \cdot 10 ^3}{12 \cdot 4 \cdot 8{,}31} =301 \,\text{К} .\]

Ответ: $301$ К.
}

% Source 7165: https://phys-ege.sdamgia.ru/problem?id=7165
\BankTask{7165}{
\begin{center}\begin{tikzpicture}[task graph,>=stealth]\draw(0,1.45)circle(1.45);\draw(-1.3,.82)--(-.48,-1.18)--(.48,-1.18)--(1.3,.82);\draw(-.58,-.91)rectangle(.58,-1.18);\foreach \x in {-.5,-.34,...,.46}{\draw(\x,-.91)--++(.06,-.27);}\draw(-.09,-.91)--(-.09,-.14)..controls(-.7,.15)and(-.4,.34)..(-.56,.57)..controls(-.71,.91)and(-.18,.97)..(-.23,1.2)..controls(-.3,1.43)and(.27,1.64)..(.3,1.28)..controls(.57,1.11)and(.38,.83)..(.51,.65)..controls(.74,.35)and(.32,.3)..(.09,-.12)--(.09,-.91);\node at(.74,.9){\(t\)};\node at(-1.31,-.26){\(t_0\)};\end{tikzpicture}\end{center}
Воздушный шар, оболочка которого имеет массу M  =  $145\,\text{кг}$ и объем V  =  $230\,\text{м}^{3}$, наполняется горячим воздухом при нормальном атмосферном давлении и температуре окружающего воздуха $t_{0}$  =  $0$ °С. Какую минимальную температуру t должен иметь воздух внутри оболочки, чтобы шар начал подниматься? Оболочка шара нерастяжима и имеет в нижней части небольшое отверстие.
}{
\smallskip\noindent\textbf{Решение}\par
На шар действуют: сила Архимеда $F_A=\rho_{\,\text{в}}gV,$ сила тяжести, действующая на воздух внутри, $mg$ и сила тяжести, действующая на оболочку, $Mg.$ По второму закону Ньютона:

\[\vec F_A + m \vec{g} + M \vec{g}=0.\]

В  проекциях на вертикаль: $F_A=mg + Mg.$
Условие подъема шара: $F_{\,\text{Арх}}. > \,\text{или} \,\text{равно} Mg + mg,$ где M  — масса оболочки, $m=\rho V$— масса воздуха внутри оболочки, отсюда:

$\rho_0gV > \,\text{или} \,\text{равно} Mg + \rho gV,$
$\rho_0V > \,\text{или} \,\text{равно} M + \rho V,$

где \ensuremath{\rho}$_{0}$  — плотность окружающего воздуха, \ensuremath{\rho} — плотность воздуха внутри оболочки, V  — объем шара.
Для воздуха внутри шара находим: $\frac{pV}{T} = \frac{m}{\mu} R,$ или $\frac{m}{V} = \frac{p \cdot \mu}{R \cdot T} = \rho,$ где p  — атмосферное давление, $\mu$  — молярная масса воздуха, T  — температура воздуха внутри шара. Соответственно, имеем плотность воздуха снаружи: $\rho_0= \frac{\mu p}{RT_0} ,$ где $T_{0}$  — температура окружающего воздуха.

$\frac{p \cdot \mu \cdot V}{R \cdot T_0} > \,\text{или} \,\text{равно} M + \frac{p \cdot \mu \cdot V}{R \cdot T} .$
$\frac{p \cdot \mu \cdot V}{R \cdot T_{\min}} = \frac{p \cdot \mu \cdot V}{R \cdot T_0} - M.$
$\frac{1}{T_{\min}} = \frac{1}{T_0} - \frac{M \cdot R}{p \cdot \mu \cdot V} ,$

$T_{\min}=T_0 \frac{p\mu V}{p\mu V - MRT_0} =273 \cdot \frac{10 ^{5} \cdot 29 \cdot 10 ^{ - 3 } \cdot 230}{10 ^{5} \cdot 29 \cdot 10 ^{ - 3 } \cdot 230 - 145 \cdot 8{,}31 \cdot 273} \approx 538\,\text{К} = 265 ^\circ \,\text{С}.$

Ответ: $265$.
}

% Source 42652: https://phys-ege.sdamgia.ru/problem?id=42652
\BankTask{42652}{
\begin{center}\begin{tikzpicture}[task graph,>=stealth]\draw[->](0,0)--(5,0)node[below]{\(V\)};\draw[->](0,0)--(0,3.7)node[left]{\(p\)};\node[below left]at(0,0){\(0\)};\fill(1.2,0.35)circle(1.25pt)node[right]{\(1\)};\fill(1.2,1.6)circle(1.25pt)node[above left]{\(2\)};\fill(3.7,3.15)circle(1.25pt)node[above right]{\(3\)};\draw(1.2,0.35)--(1.2,1.6);\draw[->](1.2,0.35)--(1.200,1.125);\draw(1.2,1.6)--(3.7,3.15);\draw[->](1.2,1.6)--(2.750,2.561);\draw(3.7,3.15)--(1.2,0.35);\draw[->](3.7,3.15)--(2.150,1.414);\foreach \y/\i/\x in {.35/1/1.2,1.6/2/1.2,3.15/3/3.7}{\draw[dashed](0,\y)--(\x,\y);\node[left]at(0,\y){\(p_{\i}\)};}\draw[dashed](1.2,0)--(1.2,.35)(3.7,0)--(3.7,3.15);\node[below]at(1.2,0){\(V_1\)};\node[below]at(3.7,0){\(V_3\)};\end{tikzpicture}\end{center}
На рисунке в координатах p–V представлен циклический процесс, проводимый с идеальным одноатомным газом. Давление идеального одноатомного газа изохорно увеличивают в $4$ раза, затем объём газа увеличивают в $2{,}5$ раза так, что давление линейно зависит от объёма и возрастает в $2$ раза. После этого газ возвращают в исходное состояние в процессе, в котором давление линейно зависит от объёма. Масса газа постоянна. Определите коэффициент полезного действия теплового двигателя, работающего по этому циклу.
}{
\smallskip\noindent\textbf{Решение}\par
Проанализируем процессы, происходящие с идеальным газом. На участке $1$−$2$ давление газа изохорно увеличивается. Следовательно, согласно закону Шарля $\frac{p}{T} =const$ его температура также увеличивается. Тогда газ не совершает работу, его внутренняя энергия увеличивается. По первому закону термодинамики $Q=A + \Delta U > 0,$ то есть газ теплоту получает.
На участке $2$−$3$ и объем, и давление увеличиваются прямо пропорционально, следовательно, температура увеличивается. Таким образом, работа газа положительна, внутренняя энергия увеличивается, газ теплоту получает.
На участке $3$−$1$ объем газа уменьшается в $2{,}5$ раза, давление уменьшается в $8$ раз, следовательно, температура уменьшается. Значит, работа газа отрицательна, внутренняя энергия уменьшается, газ теплоту отдает.
В результате анализа процессов делаем вывод, что газ получает теплоту на участках $1$−$2$ и $2$−$3$. Найдем это количество теплоты.
Для участка $1$−$2$ по первому закону термодинамики 

\[Q_{12}=\Delta U_{12}= \frac{3}{2} ( p_2 - p_1 ) V_1= \frac{3}{2} ( 4p_1 - p_1 ) V_1=4{,}5p_1V_1.\]

Найдем работу, изменение внутренней энергии и полученное количество теплоты на участке $2$−$3$:

$A_{23}= \frac{p_2 + p_3}{2} ( V_3 - V_1 ) = \frac{4p_1 + 8p_1}{2} ( 2{,}5V_1 - V_1 ) =9p_1V_1;$
$\Delta U_{23}=U_3 - U_2= \frac{3}{2} ( p_3V_3 - p_2V_2 ) = \frac{3}{2} ( 8p_1 \cdot 2{,}5V_1 - 4p_1V_1 ) =24p_1V_1;$
$Q_{23}=\Delta U_{23} + A_{23}=9p_1V_1 + 24p_1V_1=33p_1V_1.$

Таким образом, за цикл идеальный газ получил количество теплоты

\[Q=Q_{12} + Q_{23}=4{,}5p_1V_1 + 33p_1V_1=37{,}5p_1V_1.\]

Работу газа за цикл найдем как площадь треугольника под графиком:

\[A= \frac{( p_2 - p_1 ) ( V_3 - V_1 )}{2} = \frac{( 4p_1 - p_1 ) ( 2{,}5V_1 - V_1 )}{2} =2{,}25p_1V_1.\]

Тогда КПД цикла

\[\eta= \frac{A}{Q} \cdot 100\%= \frac{2{,}25p_1V_1}{37{,}5p_1V_1} \cdot 100\% = 6\%.\]

Ответ: $6$\%.
}

% Source 25414: https://phys-ege.sdamgia.ru/problem?id=25414
\BankTask{25414}{
В закрытом сосуде при температуре $100 ^\circ$ находится влажный воздух с относительной влажностью воздуха $60$\% под давлением $100$ кПа. Объем сосуда изотермически уменьшили в $2{,}5$ раза. Во сколько раз нужно увеличить абсолютную температуру воздуха без изменения объема сосуда, чтобы получить такое же конечное давление? Объемом сконденсировавшейся воды пренебречь.
}{
\smallskip\noindent\textbf{Решение}\par
Влажный воздух представляет собой смесь воздуха и водяных паров. Так как температура воздуха $100 ^\circ\,\text{С} ,$ то давление насыщенного водяного пара равно $100$ кПа, в данном случае  — давлению влажного воздуха р.
По закону Дальтона давление газа в сосуде равно сумме парциальных давлений: $p=p_1\,\text{возд} + p_{\,\text{пара}}.$
Из формулы относительной влажности воздуха следует, что давление пара $p_{\,\text{пара}}=\varphi p_{\,\text{н}}. \,\text{п}..$ Откуда давление сухого воздуха в сосуде $p_1\,\text{возд}=p ( 1 - \varphi ) .$
После изотермического сжатия водяной пар стал насыщенным, его давление стало равным давлению насыщенного пара. Давление сухого воздуха при изотермическом сжатии в $2{,}5$ раза увеличилось в $2{,}5$ раза по закону Бойля-Мариотта $pV=const,$ $p_2\,\text{возд}=2{,}5p_1\,\text{возд}.$ Тогда давление влажного воздуха в сосуде стало $p_2=p_2\,\text{возд} + p_{\,\text{н}}. \,\text{п}..$ Таким образом, $p_2=2{,}5 ( 1 - \varphi ) p + p= ( 3{,}5 - 2{,}5\varphi ) p.$
При изохорном нагревании по закону Шарля $\frac{p}{T} =const.$ Откуда:

\[\frac{T_2}{T_1} = \frac{p_2}{p} =3{,}5 - 2{,}5\varphi=3{,}5 - 2{,}5 \cdot 0{,}6=2.\]

Значит, абсолютную температуру воздуха нужно увеличить в $2$ раза.

Ответ: в $2$ раза.
}

% Source 7875: https://phys-ege.sdamgia.ru/problem?id=7875
\BankTask{7875}{
\begin{center}\begin{tikzpicture}[task graph,>=stealth]\draw(5,0)--(0,0)--(0,2.65)--(5,2.65);\draw[fill=orange!25](2.26,0)rectangle(2.66,2.65);\foreach \y in {.13,.3,...,2.63}{\draw(2.26,\y)--++(.4,-.12);}\draw(0,1.26)--(.18,1.26);\foreach \x in {.2,.45,...,1.95}{\draw (\x,1.26)arc(180:-145:.13 and .19);}\draw(2.05,1.26)--(2.26,1.26);\node at(1.05,1.91){\(k\)};\node at(.96,.55){\(p_0\)};\node at(4,1.15){\(p_0\)};\draw(2.63,2.1)--(3.8,3.14)node[right]{\(S\)};\draw(0,0)--(0,-.95)(2.26,0)--(2.26,-.95);\draw( .22,-.74)--(2.04,-.74)node[midway,above]{\(L\)};\end{tikzpicture}\end{center}
В горизонтальном цилиндре с гладкими стенками под массивным поршнем с площадью S находится одноатомный идеальный газ. Поршень соединен с основанием цилиндра пружиной с жесткостью k. В начальном состоянии расстояние между поршнем и основанием цилиндра равно L, а давление газа в цилиндре равно внешнему атмосферному давлению $p_{0}$ (см. рис.). Какое количество теплоты Q передано затем газу, если в результате поршень медленно переместился вправо на расстояние b?
}{
\smallskip\noindent\textbf{Решение}\par
Тепло, переданное газу, идет на изменение его внутренней энергии и на совершением им работы:

\[Q=\Delta U + A.\]

В начальном состоянии давление и объем газа равны $p_0$ и $V_0=SL.$ Следовательно, пружина была не деформирована. 
В конечном состоянии пружина была деформирована на величину b, возникла сила упругости $F_{\,\text{упр}}=kb,$ направленная влево. Запишем условие равновесия поршня в конечном состоянии:

\[p_0S + F_{\,\text{упр}}=p_{\,\text{к}}S.\]

Тогда давление воздуха в конечном состоянии $p_{\,\text{к}}=p_0 + \frac{kb}{S}$ и объем воздуха в конечном состоянии  $V_{\,\text{к}}=S ( L + b ) .$ 
Изменение внутренней энергии воздуха $\Delta U = \frac{3}{2} \nu R\Delta T.$ Тогда получаем:

\[\Delta U = \frac{3}{2} ( p_{\,\text{к}}V_{\,\text{к}} - p_0V_0 ) = \frac{3}{2} ( p_0Sb + kbL + kb ^2 ) .\]

Чтобы рассчитать работу, заметим, что в каждый момент времени, когда поршень сдвинут на x от начального положения давление равно $p=p_0 + \frac{kx}{S} =p_0 + \frac{k ( V - V_0 )}{S ^2} ,$ то есть давление линейно зависит от объема. Значит, на pV-диаграмме процесс расширения будет изображен отрезком прямой, а фигура под графиком будет являться трапецией, площадь которой равна:

\[A = \frac{p_0 + p_{\,\text{к}}}{2} ( V_{\,\text{к}} - V_0 ) =p_0Sb + \frac{kb ^2}{2} .\]

Таким образом, полученное воздухом количество теплоты равно:

\[Q = \frac{5}{2} p_0Sb + \frac{3}{2} kbL + 2kb ^2 .\]

Ответ: $Q = \frac{5}{2} p_0Sb + \frac{3}{2} kbL + 2kb ^2 .$
}

% Source 7129: https://phys-ege.sdamgia.ru/problem?id=7129
\BankTask{7129}{
\begin{center}\begin{tikzpicture}[task graph,>=stealth]\draw(4.7,0)--(0,0)--(0,2.65)--(4.7,2.65);\draw[fill=orange!35](2.5,0)rectangle(2.94,2.65);\draw[dashed](3.95,0)rectangle(4.39,2.65);\node at(1.3,1.62){\(p\)};\draw(2.82,1.02)--(3.57,1.94)node[above]{\(S\)};\draw(0,0)--(0,-.65)(2.5,0)--(2.5,-.91)(3.95,0)--(3.95,-.91);\draw[<->](0,-.49)--(2.5,-.49)node[midway,below]{\(L\)};\draw[<->](2.5,-.82)--(3.95,-.82)node[midway,below]{\(x\)};\end{tikzpicture}\end{center}
В горизонтальном цилиндрическом сосуде, закрытом поршнем, находится одноатомный идеальный газ. Первоначальное давление газа p  =  $4$ · $10^{5}$ Па. Расстояние от дна сосуда до поршня равно L. Площадь поперечного сечения поршня S  =  $25\,\text{см}^{2}$. В результате медленного нагревания газ получил количество теплоты Q  =  $1{,}65\,\text{кДж}$, а поршень сдвинулся на расстояние x  =  $10\,\text{см}$. При движении поршня на него со стороны стенок сосуда действует сила трения величиной $F_{\,\text{тр}}$  =  $3$ · $10^{3}$ Н. Найдите L. Считать, что сосуд находится в вакууме.
}{
\smallskip\noindent\textbf{Решение}\par
$1$.  Поршень будет медленно двигаться, если сила давления газа на поршень и сила трения со стороны стенок сосуда уравновесят друг друга: $p_{2}$S = $F_{\,\text{тр}}$, откуда $p_2= \frac{F_{\,\text{тр}}}{S} = 12 \cdot 10 ^{5} \,\text{Па} > p_1.$ $( p_1 = p ) .$
$2$.  Поэтому при нагревании газа поршень будет неподвижен, пока давление газа не достигнет значения $p_{2}$. В этом процессе газ получает количество теплоты $Q_{12}$. Затем поршень будет сдвигаться, увеличивая объем газа, при постоянном давлении. В этом процессе газ получает количество теплоты $Q_{23}$.
$3$.  В процессе нагревания, в соответствии с первым началом термодинамики, газ получит количество теплоты:

\[Q=Q_{12} + Q_{23} = ( U_3 - U_1 ) + p_2Sx = ( U_3 - U_1 ) + F_{\,\text{тр}}x.\]

$4$.  Внутренняя энергия одноатомного идеального газа: 

\[U_1= \frac{3}{2} \nu RT_1= \frac{3}{2} p_1SL\]

в начальном состоянии, 

\[U_3= \frac{3}{2} \nu RT_3= \frac{3}{2} p_2S ( L + x ) = \frac{3}{2} F_{\,\text{тр}} ( L + x )\]

в конечном состоянии.
$5$.  Из пп. $3$, $4$ получаем $L= \frac{Q - \frac{5}{2} F_{\,\text{тр}}x}{\frac{3}{2} ( F_{\,\text{тр}} - p_1S )} = \frac{1650 - 2{,}5 \cdot 3000 \cdot 0{,}1}{1{,}5 \cdot ( 3000 - 400000 \cdot 0{,}0025 )} =0{,}3\,\text{м}.$

Ответ: L  =  $0{,}3\,\text{м}$.
}

% Source 29463: https://phys-ege.sdamgia.ru/problem?id=29463
\BankTask{29463}{
\begin{center}\begin{tikzpicture}[task graph,>=stealth]\fill[cyan!22](0,2.95)--(0,.54)arc(180:270:.54)--(3.85,0)arc(270:360:.54)--(4.39,2.95)--cycle;\draw(0,3.5)--(0,.54)arc(180:270:.54)--(3.85,0)arc(270:360:.54)--(4.39,3.5);\draw(0,2.95)--(4.39,2.95);\draw[fill=white](1.02,1.43)rectangle(3.08,2.73);\foreach \x in {1.13,1.34,...,2.86}{\foreach \y in {1.56,1.78,...,2.63}{\fill[gray](\x,\y)circle(.24pt);}}\draw(3.08,2.73)--(3.68,2.73)(3.08,1.43)--(3.68,1.43);\draw[fill=orange!30](3.14,1.43)rectangle(3.39,2.73);\draw(1.12,1.23)circle(.1)(3.16,1.23)circle(.1);\draw(1.12,1.43)--(1.12,1.33)(3.16,1.43)--(3.16,1.33)(1.18,1.15)--(2.16,.14)--(3.1,1.15)(2.16,.14)--(2.16,0);\foreach \x/\y in {.31/2.63,.32/1.78,.34/1.15,.53/.34,3.6/2.52,3.72/1.83,3.69/.98,2.8/.35}{\draw[cyan!55!black](\x,\y)--++(.34,0);}\end{tikzpicture}\end{center}
Сосуд с воздухом, отделенным от атмосферы поршнем, поместили в сосуд с водой и прикрепили ко дну кастрюли нитью, не дающей сосуду всплыть на поверхность. Начальная температура воздуха в сосуде и воды в кастрюле $t_{1}$  =  $18$ °C. Кастрюлю помещают в холодильник. Ниже какой температуры $t_{2}$ должна охладиться вода, чтобы сосуд опустился на дно кастрюли?
Начальный объем воздуха в сосуде $V_{1}$  =  $2$ дм$^{3}$. Масса сосуда с поршнем m  =  $2{,}2\,\text{кг}$. Сосуд и поршень изготовлены из стали, плотность которой $\rho = 7800 \,\text{кг/м} ^3 .$ Поршень может скользить без трения со стенками сосуда. Считать массу воздуха в сосуде намного меньшей, чем масса сосуда с поршнем.
}{
\smallskip\noindent\textbf{Решение}\par
$1$.  В сосуде с воздухом по условию поршень может скользить без трения. Значит, он будет перемещаться до тех пор, пока внешнее и внутреннее давление не станут равными. Учитывая, что глубина кастрюли небольшая, внешнее давление на поршень существенно не изменится. Поэтому процесс охлаждения воздуха в сосуде можно считать изобарным. Тогда по закону Гей-Люссака:

\[\frac{V_1}{V_2} = \frac{T_1}{T_2} ,\]
 откуда:

\[T_2= \frac{V_2}{T_1} V_1.\]

$2$.  При начальной температуре воды на сосуд действуют сила Архимеда, направленная вверх, сила тяжести сосуда, направленная вниз, и результирующая сил натяжения нитей, направленная вниз. С учетом того, что масса воздуха намного меньше массы сосуда с поршнем, силой тяжести, действующей на воздух можно пренебречь. 
$3$.  При понижении температуры воды, воздух в сосуде будет изобарно охлаждаться, объем воздуха уменьшится. Тогда уменьшится суммарный объем воздуха и сосуда. Силы натяжения нитей станут равными нулю. Тело будет тонуть при условии, что $F_A < mg.$ Сила Архимеда, действующая на тело, равна $F_A=\rho _0gV,$ где $V=V_c + V_2.$ Объем сосуда с поршнем $V_c= \frac{m}{\rho} .$ 
$4$.  Объединяя уравнения и неравенство, получаем:

\[\rho _0 ( \frac{m}{\rho} + V_2 ) < m,\]
 откуда:

\[V_2 < \frac{m ( \rho - \rho _0 )}{\rho\rho _0} .\]

Следовательно, температура воды должна быть:

$T_2 < \frac{m ( \rho - \rho _0 ) T_1}{\rho \rho _0V_1} = \frac{2{,}2\,\text{кг} \cdot ( 7{,}8 - 1 ) \cdot 10 ^3 \,\text{кг/м} ^3 \cdot 291\,\text{К}}{10 ^3 \,\text{кг/м} ^3 \cdot 7{,}8 \cdot 10 ^3 \,\text{кг/м} ^3 \cdot 2 \cdot 10 ^{ - 3 } \,\text{м} ^3} \approx 279\,\text{К}$ или $t_2 < 6 ^\circ C.$

Ответ: ниже $6 ^\circ C.$
}

% Source 7935: https://phys-ege.sdamgia.ru/problem?id=7935
\BankTask{7935}{
Два одинаковых теплоизолированных сосуда соединены короткой трубкой с краном. Объем каждого сосуда V  =  $1\,\text{м}^{3}$. В первом сосуде находится $\nu_{1}$  =  $1$ моль гелия при температуре $T_{1}$  =  $400$ К; во втором  — $\nu_{2}$  =  $3$ моль аргона при температуре $T_{2}$. Кран открывают. После установления равновесного состояния давление в сосудах р  =  $5{,}4$ кПа. Определите первоначальную температуру аргона $T_{2}$.
}{
\smallskip\noindent\textbf{Решение}\par
$1$.  Поскольку в указанном процессе газ не совершает работы и система является теплоизолированной, то в соответствии с первым законом термодинамики суммарная внутренняя энергия газов сохраняется:

\[\frac{3}{2} \nu_1RT_1 + \frac{3}{2} \nu_2RT_2= \frac{3}{2} ( \nu_1 + \nu_2 ) RT,\]

где T  — температура в объединенном сосуде в равновесном состоянии после открытия крана.
$2$.  В соответствии с уравнением Клапейрона  — Менделеева для конечного состояния можно записать: 

\[p ( 2V ) = ( \nu_1 + \nu_2 ) RT.\]

Исключая из двух записанных уравнений конечную температуру T, получаем искомое выражение для начальной температуры аргона:

\[T_2= \frac{2Vp}{\nu_2R} - \frac{\nu_1}{\nu_2} T_1= \frac{2 \cdot 1 \cdot 5{,}4 \cdot 10 ^3}{3 \cdot 8{,}31} - \frac{1}{3} \cdot 400 \approx 300\,\text{К}.\]

Ответ: $T_2 \approx 300\,\text{К}.$
}

% Source 9976: https://phys-ege.sdamgia.ru/problem?id=9976
\BankTask{9976}{
Объем $0{,}1$ литра водорода нагревают при постоянном давлении от $300$ до $3000$ К. При высоких температурах молекулы водорода распадаются на отдельные атомы. На графике показана зависимость доли распавшихся молекул от температуры. Чему равен конечный объем газа?
\begin{center}\begin{tikzpicture}[task graph,>=stealth]\draw[gray!65,xstep=.6,ystep=.6](0,0)grid(5.4,3.6);\draw[->](0,0)--(5.4,0)node[below]{\(T,\,\text{К}\)};\draw[->](0,0)--(0,3.7)node[left]{\(\alpha,\,\%\)};\node[below left]at(0,0){\(0\)};\foreach \x/\t in {1.2/1000,2.4/2000,3.6/3000}{\node[below]at(\x,0){\(\t\)};}\foreach \y/\a in {.6/10,1.2/20,1.8/30,2.4/40}{\node[left]at(0,\y){\(\a\)};}\draw[smooth] plot coordinates{(0,0)(.6,.008)(1.2,.052)(1.8,.16)(2.4,.34)(3,.69)(3.6,1.2)(4.2,1.81)(4.8,2.53)(5.4,3.3)};\end{tikzpicture}\end{center}
}{
\smallskip\noindent\textbf{Решение}\par
В первоначальном состоянии весь водород находился в молекулярном состоянии, и согласно уравнению Менделеева  — Клапейрона:

\[pV_1=\nu RT_1.\]

После нагревания согласно графику $\alpha =20\%$ молекул распадется. При распаде из одной молекулы водорода образуется два атома водорода, значит, в сосуде будет $( 1 - \alpha ) \nu$ молекулярного водорода и $2 \alpha \nu$ атомарного водорода. Общее количество вещества стане равным $( 1 - \alpha ) \nu + 2 \alpha \nu = ( 1 + \alpha ) \nu = 1{,}2\nu,$ и согласно уравнению Менделеева  — Клапейрона:

\[pV_2=1{,}2\nu RT_2.\]

Поделив второе уравнение на первое, получаем:

$\frac{V_2}{V_1} =1{,}2 \frac{T_2}{T_1} =1{,}2 \cdot \frac{3000}{300} =12,$

$V_2=12V_1=12 \cdot 0{,}1=1{,}2\,\text{л}.$

Ответ: $V_2=1{,}2\,\text{л}.$
}

\end{document}
